\documentclass[10pt,a4paper]{amsart}
\usepackage[margin=19mm]{geometry}
\usepackage{amsmath,amssymb,mathtools}
\usepackage[scr=rsfs,cal=euler]{mathalfa}
\usepackage{xfrac}
\usepackage[unicode,hidelinks,bookmarks=false]{hyperref}
\newcommand{\ie}{i.e.\ }
\newcommand{\cf}{cf.\ }
\newcommand{\resp}{respectively\ }
\newcommand{\Subsection}{\subsection}
\providecommand{\nobreakdash}{\nobreak\hbox{-}}
\setlength{\emergencystretch}{2em}
\multlinegap0pt
\usepackage{amssymb}
\usepackage{xcolor}
\usepackage{float}
\newtheorem{prop}{Proposition}
\def\bP{{\mathbb P}}
\title{Equivariant unirationality of Fano threefolds}
\author{Ivan Cheltsov, Yuri Tschinkel, Zhijia Zhang}
\date{Source: arXiv:2502.19598v2 (CC BY 4.0)}
\begin{document}
\maketitle

\noindent\emph{Source and licence.} Equivariant unirationality of Fano threefolds. Version arXiv:2502.19598v2 (CC BY 4.0), distributed by arXiv under the Creative Commons Attribution 4.0 International licence, http://creativecommons.org/licenses/by/4.0/, as recorded in the arXiv licence metadata for this exact version and shown on the abstract page https://arxiv.org/abs/2502.19598v2. Full licence notice: \texttt{licenses/arxiv-2502.19598-CC-BY-4.0.txt}.\par\smallskip

\noindent\textit{Editorial scope.} Counted unit: the article's prop prop:cubic-gen (source0.tex lines 2053--2056). The article's complete original proof of it is delivered with it (source0.tex lines 2058--2067, one \texttt{proof} environment of 10 lines). No mathematics is written, repaired or re-derived: the statement and the proof are the article's own lines, in the article's own notation, and no retained line is substituted or re-flowed.  No further source-local context is needed: every cross-reference of every retained line resolves inside the two delivered spans.  Nothing was omitted from any retained span, so the package carries no omission record.  No retained line contains a citation, so the wrapper carries no bibliography and no bibliography entry.  No retained line contains a figure environment, an image-inclusion command or drawing code, so no image is delivered and the package declares no asset dependency.  Editorial formatting only, in addition: an AMS \texttt{amsart} preamble that restates the article's own declarations and macros used by the retained lines, namely the environments \texttt{prop} on separate counters, together with the standard packages those lines need.  The retained environments print in this document as Proposition 1 in order of appearance, whereas the article prints the same items with the numbering of its own full document.  The complete native source of the article as distributed in the arXiv e-print for this exact version is bundled unchanged as \texttt{sources/173927/source0.tex}.
\begin{prop}
\label{prop:cubic-gen}
Let $X\subset \bP^n$, $n\ge 4$, be a smooth cubic hypersurface with a regular, generically free action of a finite group $G$. Assume that $G$ contains a subgroup $G'$ of index 2 with isolated fixed points. Then $X$ is $G$-unirational. 
\end{prop}

\begin{proof} 
It suffices to consider the case when $G$ has no fixed points on $X$. Since $G'$ is an index-2 normal subgroup of $G$, we have (at least) one $G$-orbit of $G'$-fixed points $s_1,s_2 \in X$.
The tangent hyperplane sections $S_1,S_2$ at these points 
are irreducible cubic hypersurfaces of dimension $\ge 2$, and not cones (since $X$ is smooth). The surfaces $S_1,S_2$ are $G'$-lineariazable, 
via projection from $s_1, s_2$, as they have a distinguished $G'$-fixed singular point.  The $G$-rational map 
$$
S_1\times S_2\dashrightarrow X
$$
is dominant, which shows $G$-unirationality. 
\end{proof}

\end{document}
